\section{Hrushovski--Weil for families of bijections}
\label{sec:bij}

Let $V,V'$ be irreducible $k$-varieties.  A \emph{generic set of invertible germs
from $V$ to $V'$} is a nonempty $\Gamma$-orbit $\mathcal F\subseteq\Germ^\times(V,V')$
(invertible germs $V\to V'$, \cref{lem:inverse}); its members have a common rank
$n=\rk(\mathcal F)$ and $\dim V=\dim V'=:m$.  It is \emph{generically transitive} if
$f(y)\ind{}y$ for all $f\in\mathcal F$ and $y\in V_{\Gen}$ with $y\ind{}f$
(exactly as \ref{H:trans}).  \Cref{lem:genset-basic} holds verbatim for $\mathcal F$.

\begin{definition}[the difference set]\label{def:diff}
For $f_1,f_2\in\mathcal F$ with $f_1\ind{}f_2$ put $\delta(f_1,f_2)=f_1^{-1}f_2\in\Germ^\times(V)$,
and let $\mathcal G(\mathcal F)$ be the set of all such germs.
\end{definition}

\begin{lemma}\label{lem:diff-basic}
\begin{enumerate}[nosep]
\item $\mathcal G(\mathcal F)$ is a single $\Gamma$-orbit, hence a generic set of germs on $V$.
\item For $f_1\ind{}f_2$ in $\mathcal F$ and $\delta=\delta(f_1,f_2)$: $\delta$ and $f_2$ are
  interalgebraic over $f_1$, $\delta$ and $f_1$ are interalgebraic over $f_2$, and
  $\rk(\delta)\ge n$.  The following are equivalent: $\rk(\delta)=n$; $\delta\ind{}f_1$;
  $\delta\ind{}f_2$.
\end{enumerate}
\end{lemma}
\begin{proof}
(1) Independent pairs $(f_1,f_2)$ correspond to independent pairs of generic points of
the code variety $W$, which are all $\Gamma$-conjugate (\cref{lem:indep-conj}); and
$\sigma(f_1^{-1}f_2)=(\sigma f_1)^{-1}\sigma f_2$.
(2) $\Cb(\delta)\subseteq\acl(f_1,f_2)$, $\Cb(f_2)\subseteq\acl(f_1,\delta)$ (as $f_2=f_1\delta$),
$\Cb(f_1)\subseteq\acl(f_2,\delta)$ (as $f_1=f_2\delta^{-1}$), by \cref{lem:comp-codes}(1)
applied to $\delta:V\to V$ and $f_1:V\to V'$.  Hence
$\rk(\delta/f_1)=\rk(f_2/f_1)=n$ and $\rk(\delta)\ge n$; and $\rk(\delta)=n$ iff
$\rk(\delta)=\rk(\delta/f_1)$ iff $\delta\ind{}f_1$.  Symmetrically for $f_2$.
\end{proof}

\begin{theorem}[Hrushovski--Weil for bijections]\label{thm:hw-bij}
Let $\mathcal F$ be a generically transitive generic set of invertible germs from $V$
to $V'$ of rank $n$, and suppose
\begin{enumerate}[label=\textup{(F)}]
\item\label{F:rank} $\rk(\delta)=n$ for one (equivalently every) $\delta\in\mathcal G(\mathcal F)$.
\end{enumerate}
Then $\mathcal G:=\mathcal G(\mathcal F)$ is a generically transitive generic set of germs on
$V$ of rank $n$ closed under inverse and generic composition, so \cref{thm:hw} applies to
it.  Moreover, for $f\in\mathcal F$ and $\gamma$ in the group $G$ generated by $\mathcal G$
with $\gamma\ind{}f$: $f\gamma\in\mathcal F$ and $f\gamma\ind{}\gamma$.
\end{theorem}

\begin{proof}
\emph{Inverse \ref{H:inv}.}  $\delta(f_1,f_2)^{-1}=f_2^{-1}f_1=\delta(f_2,f_1)$.

\emph{Generic composition \ref{H:comp}.}  Let $\delta,\delta'\in\mathcal G$ with
$\delta\ind{}\delta'$, and choose $f\in\mathcal F$ with $f\ind{}\delta\delta'$
(\cref{lem:genset-basic}(1)).  Write $\delta=f_1^{-1}f_2$ with $f_1\ind{}f_2$.
By \cref{lem:diff-basic}(2) and \ref{F:rank}, $(f_2,\delta)$ is an independent pair
consisting of a member of $\mathcal F$ and a member of $\mathcal G$; so is $(f,\delta)$.
By \cref{lem:indep-conj} (applied to the code varieties of $\mathcal F$ and $\mathcal G$)
there is $\sigma\in\Gamma$ with $\sigma(f_2,\delta)=(f,\delta)$.  Put $f'_1=\sigma f_1\in\mathcal F$;
then $f'^{-1}_1f=\sigma(f_1^{-1}f_2)=\sigma\delta=\delta$.  Similarly $(f_1,\delta)$ and
$(f,\delta')$ are independent $\mathcal F\times\mathcal G$ pairs, so there is $\sigma'$ with
$\sigma'(f_1,\delta)=(f,\delta')$; put $f'_2=\sigma'f_2$, so $f^{-1}f'_2=\delta'$.  Hence
\[
   \delta\delta'=f'^{-1}_1\,f\,f^{-1}\,f'_2=f'^{-1}_1f'_2 .
\]
It remains to check that $f'_1\ind{}f'_2$.  First, $f'_1\ind{}f$ and $f'_2\ind{}f$ (images of
independent pairs).  Next, $f'_1$ and $\delta$ are interalgebraic over $f$
($\delta=f'^{-1}_1f$, $f'_1=f\delta^{-1}$), and $f'_2$ and $\delta'$ are interalgebraic
over $f$; and $\delta\ind{f}\delta'$ because $\delta\ind{}\delta'$ and $f\ind{}\delta\delta'$
(\ref{R:trans}: $\rk(\delta/f\delta')=\rk(\delta/\delta')=\rk(\delta)=\rk(\delta/f)$).
Hence $f'_1\ind{f}f'_2$ by \ref{R:aclinv}, and with $f'_1\ind{}f$ we get $f'_1\ind{}ff'_2$
(\ref{R:trans}), so $f'_1\ind{}f'_2$.  Therefore $\delta\delta'=\delta(f'_1,f'_2)\in\mathcal G$.
(That $\delta\delta'$ is then independent from $\delta$ and from $\delta'$ is
\cref{lem:prod-indep}.)

\emph{Generic transitivity \ref{H:trans}.}  Let $\delta\in\mathcal G$ and $y\in V_{\Gen}$
with $y\ind{}\delta$.  The set of pairs $(f_1,f_2)$ with $f_1\ind{}f_2$ and
$f_1^{-1}f_2=\delta$ is $\Gamma_{\acl(\Cb\delta)}$-invariant (\cref{lem:germ-aut}) and
nonempty, so by \cref{lem:fresh-conj} we may choose one with $f_1f_2\ind{\delta}y$; then
$y\ind{}\delta f_1f_2$, in particular $y\ind{}f_2$ and $y\ind{}f_1f_2$.  Put $c=f_2(y)$;
by generic transitivity of $\mathcal F$, $c\ind{}y$, and $c\in\acl(y,f_2)$.  From
$y\ind{}f_1f_2$ and $f_1\ind{}f_2$ we get $\rk(f_1/f_2y)=\rk(f_1f_2y)-\rk(f_2y)=
(m+2n)-(m+n)=n$, so $f_1\ind{}f_2y$ and hence $f_1\ind{}yc$.  Then $c\ind{}y$ and
$c\ind{y}f_1$ (from $f_1\ind{}yc$) give $c\ind{}yf_1$ (\ref{R:trans}), so $c$ is a generic
point of $V'$ over $F:=\acl(\Cb f_1,y)$.  By \cref{lem:eval}(2) applied to $f_1^{-1}$ and
$F$, $z:=f_1^{-1}(c)$ is a generic point of $V$ over $F$, so $z\ind{}y$.  Finally
$z=f_1^{-1}(f_2(y))=\delta(y)$ by \cref{lem:comp}, as $y$ is generic over
$\acl(\Cb f_1\,\Cb f_2)$.

\emph{The last claim.}  First let $f\in\mathcal F$ and $\delta\in\mathcal G$ with
$\delta\ind{}f$.  Fix any independent $f_1\ind{}f_2$ in $\mathcal F$ and $\delta_0=f_1^{-1}f_2$;
by \ref{F:rank} and \cref{lem:diff-basic}(2), $(f_1,\delta_0)$ is an independent
$\mathcal F\times\mathcal G$ pair, as is $(f,\delta)$, so \cref{lem:indep-conj} gives
$\sigma\in\Gamma$ with $\sigma(f_1,\delta_0)=(f,\delta)$.  Then $f\delta=\sigma(f_1\delta_0)=\sigma f_2\in\mathcal F$,
$f\delta\ind{}f$ (image of $f_2\ind{}f_1$) and $f\delta\ind{}\delta$ (image of
$f_2\ind{}\delta_0$, which is \ref{F:rank}).  Now let $\gamma\in G$ with $\gamma\ind{}f$.  By
\cref{thm:hw}\ref{hw:two} write $\gamma=\delta_1\delta_2$ with $\delta_i\in\mathcal G$,
conjugated over $\acl(\Cb\gamma)$ (\cref{lem:fresh-conj}) so that
$\delta_1\delta_2\ind{\gamma}f$; then
$f\ind{}\gamma\delta_1\delta_2$.  By the first part, $f_2:=f\delta_1\in\mathcal F$ with
$f_2$ and $f$ interalgebraic over $\delta_1$, so
$\rk(f_2/\delta_2)\ge\rk(f_2/\delta_1\delta_2)=\rk(f/\delta_1\delta_2)=n$, i.e.\ $f_2\ind{}\delta_2$,
and again by the first part $f\gamma=f_2\delta_2\in\mathcal F$.  Finally $f\gamma$ and $f$
are interalgebraic over $\gamma$ (\cref{lem:comp-codes}(1)), so
$\rk(f\gamma/\gamma)\ge\rk(f\gamma/\delta_1\delta_2)=\rk(f/\delta_1\delta_2)=n$ (as
$\Cb(\gamma)\subseteq\acl(\delta_1\delta_2)$), i.e.\ $f\gamma\ind{}\gamma$.
\end{proof}

\begin{remark}
In the proof of generic composition the only use of \ref{F:rank} is to know that
$(f_2,\delta)$ and $(f_1,\delta)$ are independent pairs, which is what allows
``re-basing'' the representation $\delta=f_1^{-1}f_2$ on the fresh element $f$
(it is used once more, in the last claim, for $f\delta\ind{}\delta$).  In the
finite-rank setting \ref{F:rank} is exactly the hypothesis ``$c\ind{}b_i$'' of Bays'
Lemma 5.4, by \cref{lem:diff-basic}(2).  In \cref{sec:proof} it is verified for the
configuration using the non-degeneracy clause.
\end{remark}
